\chapter{Data and Evaluation Setup}

\section{Formal Problem Definition}
To evaluate model drift rigorously, we must first define the classification task within a probabilistic framework. Let $\mathcal{X} \subset \mathbb{R}^d$ denote the high-dimensional input space of social-media posts (comprising text, metadata, and visual features) and $\mathcal{Y} = \{0, 1\}$ the binary label space, where $1$ indicates relevance to a campaign. 

In the context of LLM-as-a-Service, the model is a non-stationary black box. We define the classifier as a function $f(x; \pi, \theta_t) \to \hat{y} \in \{0, 1\}$, where $\pi$ is a fixed prompt and $\theta_t$ represents the hidden, time-varying model weights. 

Under the zero-one loss function $\mathcal{L}(\hat{y}, y) = \mathbb{I}(\hat{y} \neq y)$, the \textbf{Expected Risk} (or true error rate) at time $t$ is:
\begin{equation}
    R_t(f) = \mathbb{E}_{(x,y) \sim P_t} [\mathbb{I}(f(x; \pi, \theta_t) \neq y)]
\end{equation}
Since the true data-generating distribution $P_t$ is unknown, we must estimate this risk empirically.

\section{The Reference Benchmark vs. Operational Constraints}
Let $\mathcal{D}_N$ be the reference benchmark consisting of $N = 1199$ manually labelled examples, establishing the ground truth. Evaluating all candidate models against $\mathcal{D}_N$ yields a highly precise empirical risk $\hat{R}_N$. 

However, running inference on $N=1199$ instances weekly across multiple candidate models is economically prohibitive. We require a monitoring subset $\mathcal{D}_n \subset \mathcal{D}_N$ of size $n \ll N$. The central mathematical challenge is to select $\mathcal{D}_n$ such that it remains highly sensitive to temporal shifts in the parameters $\theta_t$, despite the significantly reduced sample size.

\section{Reference benchmark and monitoring subset}
\label{sec:input-datasets}

The two sets $\mathcal{D}_N$ and $\mathcal{D}_n$ are stored as fixed Langfuse datasets sharing one schema and label column (\texttt{campaign\_relevant}). Table~\ref{tab:dataset-summary} summarises their composition.

\begin{table}[htbp]
    \centering
    \caption{Full reference benchmark (\texttt{campaign\_relevance\_02e1a68ccb0f}) and monitoring subset (\texttt{campaign\_relevance\_disagree\_subset\_9d488308aa46}): size, class balance, and modality composition.}
    \label{tab:dataset-summary}
    \begin{tabular}{lcc}
        \toprule
        & \textbf{Full benchmark} & \textbf{Monitoring subset} \\
        \midrule
        Number of posts ($N$ / $n$) & 1199 & 300 \\
        Relevant (\texttt{campaign\_relevant}${}=1$) & 682 (56.9\%) & 250 (83.3\%) \\
        Not relevant (\texttt{campaign\_relevant}${}=0$) & 517 (43.1\%) & 50 (16.7\%) \\
        Image posts & 288 & 78 \\
        Video posts & 911 & 222 \\
        Total attached images (all posts) & 5911 & 1762 \\
        Mean images per post & 4.93 & 5.87 \\
        \bottomrule
    \end{tabular}
\end{table}

Each post is tagged as an \textbf{image post} (one image) or a \textbf{video post} (several frames). Both sets are \textbf{class-imbalanced}; the subset is far more positive-heavy than $\mathcal{D}_N$ (Table~\ref{tab:dataset-summary}), which follows from disagreement-based construction (Section~\ref{sec:informative-sampling-disagreement}). Accuracy figures on $\mathcal{D}_N$ and $\mathcal{D}_n$ are not interchangeable---prevalence differs, not only sample size.

Roughly three quarters of items are \textbf{video posts} in both splits; multi-frame requests carry more multimodal tokens than single-image ones. The subset is not a stratified copy of the full benchmark: it has a slightly larger fraction of image posts (78/300 vs.\ 288/1199) but still a \textbf{higher} mean images per post (5.87 vs.\ 4.93), so monitored runs are visually heavier and costlier to repeat. Subsequent empirical sections report metrics on $\mathcal{D}_n$ only; they characterise model behaviour on this fixed disagreement subset, not a simple random replicate of $\mathcal{D}_N$'s label prevalence or its image--video mix.

\section{Informative Sampling via Model Disagreement}
\label{sec:informative-sampling-disagreement}
A simple random sample (SRS) of size $n$ would yield an unbiased estimator of $R_t$. However, the vast majority of social-media posts are "easy" instances located far from the model's decision boundary, where the conditional probability $P(Y=1|X) \approx 1$ or $0$. These instances exhibit low prediction variance and are highly robust to small parameter shifts in $\theta_t$. 

To maximize the sensitivity of our monitoring subset, we abandon uniform sampling in favor of \textbf{Uncertainty Sampling}. We construct $\mathcal{D}_n$ exclusively from instances where an initial ensemble of distinct models $\mathcal{M} = \{f_1, f_2, \dots, f_K\}$ exhibited disagreement:
\begin{equation}
    \mathcal{D}_{n} = \{ x_i \in \mathcal{D}_{N} \mid \exists f_j, f_k \in \mathcal{M} : \hat{y}_{i,j} \neq \hat{y}_{i,k} \}
\end{equation}

The purpose of this subset is not to estimate the global benchmark accuracy as efficiently as possible. Its purpose is to improve \emph{drift sensitivity}: if provider-side changes alter the classifier, those changes should appear first and most clearly on cases that are already hard or unstable. This is a statistical rather than geometric claim.

\subsection{Why Disagreement Cases Are Informative}
For each benchmark item $x_i$, let
\begin{equation}
    r_i = \frac{1}{K}\sum_{m=1}^{K} \mathbb{I}\!\left(\hat{y}_{i,m}=1\right)
\end{equation}
denote the empirical fraction of ensemble models predicting relevance, and define the corresponding disagreement score
\begin{equation}
    d_i = 4r_i(1-r_i) \in [0,1].
\end{equation}
The score $d_i$ is zero under complete agreement ($r_i \in \{0,1\}$) and maximal at $r_i=0.5$, where the ensemble is evenly split. Thus $d_i$ is a simple observable measure of how controversial or unstable an instance is across candidate models.

This yields an operational justification for disagreement-based monitoring. If a future provider-side update changes the behavior of a deployed model, the probability that the item flips from correct to incorrect, or vice versa, is expected to be larger on items with high $d_i$ than on items where all models already agree. In other words, disagreement cases are selected not because they provably lie on a smooth decision manifold, but because they are empirically the cases on which small modeling differences are already visible.

\begin{proposition}[Disagreement Increases Sensitivity to Temporal Flips]
Let $E_{i,t}=\mathbb{I}(\hat{y}_{i,t}=y_i)$ denote correctness of the deployed model on item $i$ at time $t$, and define the flip indicator between two monitoring dates $t_1<t_2$ by
\begin{equation}
    F_i(t_1,t_2)=\mathbb{I}\!\left(E_{i,t_1}\neq E_{i,t_2}\right).
\end{equation}
If the flip probability $P(F_i(t_1,t_2)=1 \mid d_i)$ is non-decreasing in the disagreement score $d_i$, then for any threshold $\tau \in [0,1]$,
\begin{equation}
    \mathbb{E}[F_i(t_1,t_2)\mid d_i \ge \tau] \ge \mathbb{E}[F_i(t_1,t_2)\mid d_i < \tau].
\end{equation}
\end{proposition}

\begin{proof}
The claim follows immediately from the monotonicity assumption. Conditioning on the event $d_i \ge \tau$ shifts the distribution of $d_i$ toward larger values, and hence cannot decrease the conditional mean of $F_i(t_1,t_2)$ when the regression function $d \mapsto P(F_i(t_1,t_2)=1\mid d)$ is non-decreasing.
\end{proof}

The proposition is intentionally modest. It does not claim that disagreement items are an unbiased sample of the population, nor that they are literally the Bayes-optimal uncertainty set. It claims only that, under a transparent monotonicity assumption, filtering toward high-disagreement cases increases the expected rate of observable temporal flips. This is exactly the quantity relevant for a paired drift test such as McNemar's test.

\section{Statistical Derivation of the Sample Size}
Because the monitoring rule is based on a paired McNemar test, the sample-size argument must be expressed in terms of \emph{discordant pairs} rather than in terms of the variance of a single Bernoulli accuracy estimate.

Let $b$ denote the number of items that were correct at time $t_1$ and incorrect at time $t_2$, and let $c$ denote the number of items that were incorrect at time $t_1$ and correct at time $t_2$. McNemar's test uses only the discordant total
\begin{equation}
    m = b+c
\end{equation}
and, under the null hypothesis of no directional drift, treats
\begin{equation}
    B \mid m \sim \mathrm{Binomial}(m,0.5),
\end{equation}
where $B=b$ counts the regressions among the discordant pairs.

The practically relevant effect is not the marginal change in accuracy alone, but the directional imbalance
\begin{equation}
    \delta = p_b - p_c,
\end{equation}
where $p_b=P(E_{i,t_1}=1,E_{i,t_2}=0)$ and $p_c=P(E_{i,t_1}=0,E_{i,t_2}=1)$. Since
\begin{equation}
    \hat{A}_{t_2}-\hat{A}_{t_1}=\frac{c-b}{n},
\end{equation}
detecting a degradation of at least $\Delta=0.05$ in accuracy corresponds to detecting
\begin{equation}
    \delta \ge 0.05.
\end{equation}

To connect this effect size to McNemar's test, write
\begin{equation}
    \rho = p_b + p_c
\end{equation}
for the probability that an item is discordant between the two monitoring dates. Conditional on being discordant, the probability that the discordance is a regression is
\begin{equation}
    \pi = \frac{p_b}{p_b+p_c} = \frac{1}{2} + \frac{\delta}{2\rho}.
\end{equation}
Hence McNemar's test is more powerful when the monitored subset yields more discordant pairs, that is, when $\rho$ is larger.

For a one-sided test at level $\alpha$, a normal approximation to the exact binomial test gives the design requirement
\begin{equation}
    n \ge \frac{1}{\rho} \left(
        \frac{
            z_{1-\alpha}\sqrt{0.25} + z_{1-\beta}\sqrt{\pi(1-\pi)}
        }{
            \pi-0.5
        }
    \right)^2,
\end{equation}
where $1-\beta$ is the desired power. This formula makes the trade-off explicit:
\begin{itemize}
    \item if disagreement sampling successfully raises $\rho$, the same directional drift can be detected with fewer monitored items;
    \item if temporal flips are rare, McNemar's test has little power regardless of the nominal sample size $n$;
    \item therefore the effective design parameter is the expected number of discordant pairs $m \approx n\rho$, not $n$ in isolation.
\end{itemize}

To obtain an operational target, consider the practically relevant drift threshold $\Delta=0.05$ and a conventional power target of $1-\beta=0.8$ at significance level $\alpha=0.05$. On the disagreement subset, a reasonable planning regime is that roughly $10\%$ to $15\%$ of items become discordant between two monitoring dates, so $\rho \in [0.10,0.15]$. Under a $5$ percentage-point degradation, this implies
\begin{equation}
    \pi = \frac{1}{2} + \frac{0.05}{2\rho}
    \in [0.667,0.75].
\end{equation}
Substituting $(\alpha,\beta)=(0.05,0.2)$, with $z_{1-\alpha}\approx1.645$ and $z_{1-\beta}\approx0.842$, gives a required discordant count of roughly
\begin{equation}
    m = n\rho \approx 23 \text{ to } 29.
\end{equation}
This translates into a total monitoring size of roughly
\begin{equation}
    n \approx \frac{m}{\rho} \in [190,290].
\end{equation}

On this basis, setting the monitoring subset size to $n=300$ is defensible as a conservative design choice: it is large enough to deliver around $30$--$45$ discordant pairs when $\rho$ lies in the intended operating range, which is sufficient for a one-sided McNemar alert to detect a degradation of about $5$ percentage points with usable power. The number $300$ should therefore be interpreted as a practical planning value under an assumed discordance regime, not as a universal mathematical minimum.

\section{The Role of Selection Bias and Importance Weighting}

Because $\mathcal{D}_{300}$ intentionally over-represents high-disagreement cases, the empirical accuracy calculated directly on this subset, $\hat{A}_n$, is generally a \textbf{biased estimator} of the global population accuracy $A_N$. In practice, one expects these items to be harder than a uniform sample from the full benchmark, so the subset accuracy will often be lower than the full-benchmark accuracy, but the sign and magnitude of that bias are empirical rather than universal mathematical facts.

In standard statistical evaluation, selection bias is viewed as a flaw to be corrected. In the context of our drift monitoring framework, however, this bias is an operational requirement. We track the biased metric $\hat{A}_n$ because it isolates the exact region of the input space where temporal parameter shifts ($\Delta \theta_t$) manifest most strongly, thereby maximizing the statistical power of the alerting rule.

Nevertheless, it is useful to relate this heavily biased subset to the full reference benchmark, rather than treating it as an arbitrary collection of hard cases. The relevant statistical language here is \textbf{importance weighting}: Horvitz--Thompson (HT) estimation of $A_N$ is defined for sampling designs in which inclusion is random with \emph{known} probabilities $q_i=P(S_i=1)$ that are strictly positive for every benchmark item ($q_i>0$ for all $i$). Only in that setting do inverse-probability weights $1/q_i$ map a biased sample back to the uniform population measure on the benchmark.

\subsection{The Horvitz--Thompson Link}
For a benchmark of size $N$, consider the classical setup in which $S_i \in \{0, 1\}$ indicates inclusion of the $i$-th reference instance and the design specifies known inclusion probabilities $q_i=P(S_i=1)$ with $q_i>0$ for all $i=1,\dots,N$. Then the global accuracy $A_N$ can be \emph{estimated unbiasedly} from the observed sample using the \textbf{Horvitz--Thompson (HT) estimator}:
\begin{equation}
    \hat{A}_{HT} = \frac{1}{N} \sum_{i=1}^{N} \frac{S_i \cdot \mathbb{I}(\hat{y}_i = y_i)}{q_i} = \frac{1}{N} \sum_{i \in \mathcal{D}_n} \frac{\mathbb{I}(\hat{y}_i = y_i)}{q_i}
\end{equation}

Here, the inverse probability $1/q_i$ acts as an importance weight that maps the biased sampling design back to the uniform population measure on the benchmark, \emph{provided} the positivity condition holds for every unit.

\begin{proof}[Proof of Unbiasedness under a positive-probability design]
Under the stated sampling mechanism, to show that $\hat{A}_{HT}$ is unbiased for the global accuracy $A_N$, we take its expectation over the design:
\begin{equation}
    \mathbb{E} [\hat{A}_{HT}] = \mathbb{E} \left[ \frac{1}{N} \sum_{i=1}^{N} \frac{S_i \cdot \mathbb{I}(\hat{y}_i = y_i)}{q_i} \right]
\end{equation}
By the linearity of expectation, and noting that the true labels and model predictions at a fixed time $t$ are deterministic with respect to the sampling mechanism:
\begin{equation}
    \mathbb{E} [\hat{A}_{HT}] = \frac{1}{N} \sum_{i=1}^{N} \frac{\mathbb{I}(\hat{y}_i = y_i)}{q_i} \mathbb{E}[S_i]
\end{equation}
Since $S_i$ is a Bernoulli random variable, its expected value is exactly its probability of occurrence, $\mathbb{E}[S_i] = q_i$. Substituting this yields:
\begin{equation}
    \mathbb{E} [\hat{A}_{HT}] = \frac{1}{N} \sum_{i=1}^{N} \frac{\mathbb{I}(\hat{y}_i = y_i)}{q_i} \cdot q_i = \frac{1}{N} \sum_{i=1}^{N} \mathbb{I}(\hat{y}_i = y_i) = A_N
\end{equation}
\end{proof}

\subsection{Operational Utility: Signal Amplification}
The HT construction above applies to random designs with $q_i>0$ across the full benchmark. The monitoring pipeline implemented in this thesis \emph{deliberately} uses a different rule: the subset is built deterministically from model--model disagreement (with a fixed-size selection among those cases), so agreement cases are excluded by construction and their operational inclusion probability is effectively zero. Horvitz--Thompson weights $1/q_i$ are therefore undefined for the bulk of the population under this design, and the disagreement subset alone does not support reconstruction of global accuracy $A_N$. To report $A_N$ alongside the same deployment, one would add a separate probabilistic component---for example a small simple random sample (SRS) from the full benchmark in addition to the uncertainty-focused subset---so that positive, known inclusion probabilities attach to all units of interest.

For monitoring, we prioritize the high-gain signal $\hat{A}_n$: the operational loop reports raw empirical accuracy on the disagreement subset rather than an HT-style estimate, because the objective is drift sensitivity, not population-level recovery of $A_N$. The rationale is rooted in \textbf{Signal-to-Noise Ratio (SNR) optimization}. Where applicable, importance weights $1/q_i$ smooth toward the global mean---appropriate for population estimation but diluting boundary errors that matter for \textbf{anomaly detection}. By using the raw subset metric, we preserve magnification of instability where the model is least stable. Thus $\hat{A}_n$ serves as a leading indicator of drift; it is not an unbiased estimator of $A_N$, and it is not HT-based recovery of full-benchmark performance, because the implemented system is not an HT-eligible design.

